\documentclass[11pt,a4paper]{article}
\usepackage[utf8]{inputenc}
\usepackage[T1]{fontenc}
\usepackage[portuguese]{babel}
\usepackage{xcolor}
\usepackage{hyperref}
\usepackage{geometry}
\usepackage{listings}
\usepackage{tcolorbox}
\geometry{margin=2.5cm}
\definecolor{primary}{RGB}{41,128,185}
\definecolor{secondary}{RGB}{44,62,80}
\definecolor{codebg}{RGB}{240,240,240}
\lstset{backgroundcolor=\color{codebg},basicstyle=\ttfamily\small,breaklines=true,frame=single}
\title{\color{secondary}\Huge\textbf{Pipeline de Implantação e Publicação}}
\author{\color{primary}DevOps com Kubernetes -- 2026}
\date{}
\begin{document}
\maketitle


\section*{\color{primary}O que você aprenderá nesta página}
\begin{itemize}
\item Construir e publicar imagens de contêiner
\item Automatizar deploy com pipeline
\item Reduzir comandos manuais na entrega da aplicação
\end{itemize}

Publicar uma aplicação manualmente funciona nas primeiras vezes, mas rapidamente vira fonte de erro. Um pipeline de implantação organiza o caminho entre alteração de código, construção da imagem, publicação no registry e atualização dos manifestos no cluster.

\section*{\color{primary}Fluxo básico}

Um fluxo mínimo de publicação passa por quatro etapas: testar a aplicação, construir a imagem Docker, publicar a imagem em um registry e atualizar o cluster com a nova versão.

\begin{lstlisting}[language=bash]
docker build -t registry.example.com/app:v1 .
docker push registry.example.com/app:v1
kubectl set image deployment/app-dep app=registry.example.com/app:v1
\end{lstlisting}

Esse exemplo é útil para entender o processo, mas em um projeto real preferimos registrar a versão no manifesto e aplicar de forma declarativa.

\section*{\color{primary}Tags de imagem}

Usar \texttt{latest} parece conveniente, mas dificulta rastrear qual versão está em execução. Tags baseadas em commit, versão semântica ou número de build tornam o deploy reproduzível.

\begin{lstlisting}[language=bash]
containers:
  - name: app
    image: registry.example.com/app:sha-abc123
\end{lstlisting}

\section*{\color{primary}GitHub Actions como exemplo}

\begin{lstlisting}[language=bash]
name: deploy
on:
  push:
    branches: [main]
jobs:
  publish:
    runs-on: ubuntu-latest
    steps:
      - uses: actions/checkout@v4
      - name: Build image
        run: docker build -t registry.example.com/app:sha-do-commit .
      - name: Push image
        run: docker push registry.example.com/app:sha-do-commit
\end{lstlisting}

A etapa de deploy pode aplicar manifestos diretamente, mas uma evolução natural é mover isso para GitOps, como será visto na parte de operação e automação. Não coloque credenciais diretamente no arquivo do pipeline; use secrets do provedor de CI.

\end{document}
